% Options for packages loaded elsewhere
\PassOptionsToPackage{unicode}{hyperref}
\PassOptionsToPackage{hyphens}{url}
\documentclass[
  english,
]{article}
\usepackage{xcolor}
\usepackage{amsmath,amssymb}
% section numbering enabled in preamble
\usepackage{iftex}
\ifPDFTeX
  \usepackage[T1]{fontenc}
  \usepackage[utf8]{inputenc}
  \usepackage{textcomp} % provide euro and other symbols
\else % if luatex or xetex
  \usepackage{unicode-math} % this also loads fontspec
  \defaultfontfeatures{Scale=MatchLowercase}
  \defaultfontfeatures[\rmfamily]{Ligatures=TeX,Scale=1}
\fi
\usepackage{lmodern}
\ifPDFTeX\else
  % xetex/luatex font selection
\fi
% Use upquote if available, for straight quotes in verbatim environments
\IfFileExists{upquote.sty}{\usepackage{upquote}}{}
\IfFileExists{microtype.sty}{% use microtype if available
  \usepackage[]{microtype}
  \UseMicrotypeSet[protrusion]{basicmath} % disable protrusion for tt fonts
}{}
\makeatletter
\@ifundefined{KOMAClassName}{% if non-KOMA class
  \IfFileExists{parskip.sty}{%
    \usepackage{parskip}
  }{% else
    \setlength{\parindent}{0pt}
    \setlength{\parskip}{6pt plus 2pt minus 1pt}}
}{% if KOMA class
  \KOMAoptions{parskip=half}}
\makeatother
\usepackage{longtable,booktabs,array}
\newcounter{none} % for unnumbered tables
\usepackage{calc} % for calculating minipage widths
% Correct order of tables after \paragraph or \subparagraph
\usepackage{etoolbox}
\makeatletter
\patchcmd\longtable{\par}{\if@noskipsec\mbox{}\fi\par}{}{}
\makeatother
% Allow footnotes in longtable head/foot
\IfFileExists{footnotehyper.sty}{\usepackage{footnotehyper}}{\usepackage{footnote}}
\makesavenoteenv{longtable}
\usepackage{graphicx}
\makeatletter
\newsavebox\pandoc@box
\newcommand*\pandocbounded[1]{% scales image to fit in text height/width
  \sbox\pandoc@box{#1}%
  \Gscale@div\@tempa{\textheight}{\dimexpr\ht\pandoc@box+\dp\pandoc@box\relax}%
  \Gscale@div\@tempb{\linewidth}{\wd\pandoc@box}%
  \ifdim\@tempb\p@<\@tempa\p@\let\@tempa\@tempb\fi% select the smaller of both
  \ifdim\@tempa\p@<\p@\scalebox{\@tempa}{\usebox\pandoc@box}%
  \else\usebox{\pandoc@box}%
  \fi%
}
% Set default figure placement to htbp
\def\fps@figure{htbp}
\makeatother
\ifLuaTeX
\usepackage[bidi=basic,shorthands=off]{babel}
\else
\usepackage[bidi=default,shorthands=off]{babel}
\fi
\ifLuaTeX
  \usepackage{selnolig} % disable illegal ligatures
\fi
\setlength{\emergencystretch}{3em} % prevent overfull lines
\providecommand{\tightlist}{%
  \setlength{\itemsep}{0pt}\setlength{\parskip}{0pt}}
% Appended by build_report_latex.py (after pandoc header). Tectonic-compatible.

\usepackage{xcolor}
\definecolor{ReportAccent}{HTML}{1F4E79}
\definecolor{ReportMuted}{HTML}{5B6770}
\definecolor{ReportLight}{HTML}{E8F1F8}

\PassOptionsToPackage{colorlinks=true,linkcolor=ReportAccent,urlcolor=ReportAccent,citecolor=ReportAccent}{hyperref}

\usepackage{float}
\usepackage[margin=2.2cm, headheight=15pt]{geometry}
\usepackage{caption}
\usepackage{subcaption}
\usepackage[most]{tcolorbox}
\usepackage{titlesec}
\usepackage{fancyhdr}
\usepackage{enumitem}

\floatplacement{figure}{H}

\graphicspath{{figures/}}

% Numbered sections with accent colour and rule
\titleformat{\section}
  {\normalfont\Large\bfseries\color{ReportAccent}}
  {\thesection}{0.75em}{}
  [\vspace{0.15em}{\color{ReportAccent!35}\titlerule[0.6pt]}]

\titleformat{\subsection}
  {\normalfont\large\bfseries\color{ReportAccent!85!black}}
  {\thesubsection}{0.6em}{}

\titlespacing*{\section}{0pt}{1.6em}{0.8em}
\titlespacing*{\subsection}{0pt}{1.2em}{0.5em}

% Figures
\captionsetup{
  font=small,
  labelfont={bf,color=ReportAccent},
  format=plain,
  labelsep=period,
  skip=8pt,
  singlelinecheck=false,
}

% Abstract panel
\newtcolorbox{abstractbox}{
  enhanced,
  colback=ReportLight,
  colframe=ReportAccent,
  boxrule=0.8pt,
  arc=2pt,
  left=12pt,
  right=12pt,
  top=10pt,
  bottom=10pt,
  before skip=0.5em,
  after skip=1.2em,
  before upper={\setlength{\parskip}{0.85em}},
}

% Headline statistics
\newtcolorbox{keystats}{
  enhanced,
  colback=ReportLight!45,
  colframe=ReportAccent!55,
  boxrule=0.5pt,
  arc=2pt,
  left=10pt,
  right=10pt,
  top=8pt,
  bottom=8pt,
  before skip=0.6em,
  after skip=1em,
}
\newenvironment{keystatslist}{%
  \begin{itemize}[leftmargin=1.2em, itemsep=0.25em, topsep=0.2em, parsep=0pt]
}{%
  \end{itemize}
}

% Compact scientific tables
\newenvironment{scitable}{%
  \par\vspace{0.4em}
  \begin{center}
  \footnotesize
  \setlength{\tabcolsep}{4pt}
  \renewcommand{\arraystretch}{1.25}
}{%
  \end{center}
  \par\vspace{0.6em}
}

\captionsetup[table]{font=small, labelfont=bf, skip=6pt}

% Running header / footer
\pagestyle{fancy}
\fancyhf{}
\fancyhead[L]{\small\color{ReportMuted}\nouppercase{\leftmark}}
\fancyhead[R]{\small\color{ReportMuted}\thepage}
\renewcommand{\headrulewidth}{0pt}
\renewcommand{\headrule}{\color{ReportAccent!30}\hrule width\headwidth height 0.4pt}
\fancypagestyle{plain}{%
  \fancyhf{}
  \fancyfoot[C]{\small\color{ReportMuted}\thepage}
  \renewcommand{\headrulewidth}{0pt}
}

\newcommand{\um}{\,\text{$\mu$m}}
\newcommand{\umsq}{\,\text{$\mu$m}^2}

\setcounter{secnumdepth}{2}
\usepackage{bookmark}
\IfFileExists{xurl.sty}{\usepackage{xurl}}{} % add URL line breaks if available
\urlstyle{same}
\hypersetup{
  pdftitle={Frog Blood Cell Morphology --- Full Dataset Report},
  pdflang={en},
  hidelinks,
  pdfcreator={LaTeX via pandoc}}

\title{Frog Blood Cell Morphology --- Full Dataset Report}
\author{}
\date{23 May 2026}

\begin{document}
\begin{titlepage}
\thispagestyle{empty}
\centering
\vspace*{2.2cm}
{\color{ReportAccent}\rule{\textwidth}{2pt}\par}
\vspace{1.4cm}
{\Huge\bfseries\color{ReportAccent} Frog Blood Cell Morphology\par}
\vspace{0.35em}
{\Huge\bfseries\color{ReportAccent} Full Dataset Report\par}
\vspace{1.2cm}
{\large\color{ReportMuted} Morphometric analysis report\par}
\vfill
{\large 23 May 2026\par}
\vspace{2cm}
{\color{ReportAccent}\rule{\textwidth}{2pt}\par}
\end{titlepage}
\setcounter{page}{1}

\section{Abstract}\label{abstract}
\begin{abstractbox}
This report describes the \textbf{size and shape of frog blood cells}
measured from microscopy images. We focus on cells that passed an
automatic quality check and were labelled \textbf{good} for analysis.
Cells labelled bad or rejected were not included in the size summaries.

The dataset covers \textbf{479 frogs} and \textbf{7,214 images}. In
total, \textbf{383,915} detected cells were scored; \textbf{57,520}
(\textbf{15.0\%}) were accepted as good. All sizes below refer to this
good-cell set only.

\textbf{Cell size.} The typical (median) cell covers \textbf{290.69
\(\mu\)m\(^2\)} of area. Half of good cells fall between \textbf{257.70}
and \textbf{335.48 \(\mu\)m\(^2\)} (the middle 50\% of the
distribution). That area is similar to a \textbf{circle about 19.24
\(\mu\)m wide} (a simple way to picture area as a length; the cell is
not necessarily round). The average cell area was \textbf{303.10 \(\pm\)
67.07 \(\mu\)m\(^2\)} (mean \(\pm\) SD across all good cells).

\textbf{Nucleus size.} For \textbf{57,218} cells we also measured the
nucleus. The N/C ratio is nucleus area divided by cell area (a number
between 0 and 1). The typical value was \textbf{0.13}; the mean \(\pm\)
SD across cells was \textbf{0.14 \(\pm\) 0.04}.

\textbf{Differences between frogs.} Cell size was \textbf{not the same
in every frog}. Each frog has an average cell area; those averages
ranged from \textbf{212.55} to \textbf{482.23 \(\mu\)m\(^2\)}. The
\textbf{typical frog's average} (median across frogs, each frog counted
once) was \textbf{290.34 \(\mu\)m\(^2\)}. Roughly \textbf{0.661} of the
variation in cell area between cells is due to \textbf{which frog} the
cell came from. The rest is how much cells differ \textbf{within the
same frog} --- typical spread about \textbf{11.7\%} of that frog's mean
area (e.g.~\(\approx\pm\) 12\% around a frog's average). In short: frogs
differ in average cell size, but cells from one frog are usually similar
to each other.

\textbf{Nucleus vs cell size.} When cells get larger, the nucleus does
not grow as fast as the cell body. Larger cells therefore tend to have a
\textbf{proportionally smaller nucleus} (a lower N/C ratio).
\end{abstractbox}

\section{Methods}\label{methods}

\textbf{Images and cohort.} We analysed microscopy images of
\textbf{frog blood cells} from \textbf{479 frogs} (\textbf{7,214
images}). Each detected cell was linked to a \textbf{frog ID} parsed
from the image file name. All size measurements are reported in
\textbf{micrometres (\(\mu\)m)} or \textbf{square micrometres
(\(\mu\)m\(^2\))} after converting from pixels.

\textbf{Segmentation.} Cells were segmented automatically from each
image using \textbf{Cellpose} (membrane/nucleus masks). Each connected
region in the cell mask is treated as one candidate cell.

\textbf{Quality classification.} Every cell crop was scored by a trained
\textbf{ResNet18} classifier, which outputs \textbf{p(good)} --- the
model's estimated probability that the cell quality is good enough to
measure. We used \textbf{selective rejection} with three outcomes:

\begin{keystats}
\begin{itemize}
\tightlist
\item
  \textbf{p(good) \(\leq\) 0.10} \(\rightarrow\) \textbf{bad} (poor
  quality; excluded from size analysis)
\item
  \textbf{p(good) \(\geq\) 0.76} \(\rightarrow\) \textbf{good} (good
  quality; included in all morphometry and summaries)
\item
  \textbf{0.10 \textless{} p(good) \textless{} 0.76} \(\rightarrow\)
  \textbf{rejected} (model uncertain; excluded from size analysis)
\end{itemize}
\end{keystats}

Of \textbf{383,915} scored detections, \textbf{57,520} (\textbf{15.0\%})
were labelled good, \textbf{5,535} were rejected as uncertain, and the
remainder were bad. \textbf{Only good cells} enter the size results in
this report.

\textbf{Morphometry.} For each good cell we measured:

\begin{keystats}
\begin{itemize}
\tightlist
\item
  \textbf{Cell area} (\(\mu\)m\(^2\)) and \textbf{long / short diameter}
  (\(\mu\)m)
\item
  \textbf{Nucleus area and axes}, \textbf{N/C ratio} (nucleus area
  \(\div\) cell area), and \textbf{axis ratio} (major \(\div\) minor
  axis)
\end{itemize}
\end{keystats}

Pixel measurements were converted to \(\mu\)m using the \textbf{pixel
size read from each image's metadata} (OME-TIFF / TIFF), when available.

\textbf{Aggregation and yield.} For each frog we computed the
\textbf{mean and SD} of cell area (and other metrics) across its good
cells. \textbf{Yield} (good fraction) is the share of scored detections
labelled good, reported per frog and per image to show how much of each
sample passed the quality filter.

\textbf{Statistical analysis.}

\begin{keystats}
\begin{itemize}
\tightlist
\item
  \textbf{Reference intervals:} 2.5th, 50th (median), and 97.5th
  percentiles of good-cell measurements
\item
  \textbf{Between-frog variation:} ICC --- the share of cell-area
  variation explained by \textbf{which frog} a cell came from
\item
  \textbf{Within-frog spread:} for each frog, CV = (SD \(\div\) mean)
  \(\times\) 100\%; we report the median across frogs
\item
  \textbf{Nucleus--cell scaling:} log--log regression of nucleus size vs
  cell size; slope \(\approx\) 1 means nucleus and cell grow together;
  slope \textless{} 1 means the nucleus grows more slowly
\item
  \textbf{Mixed-effects models:} the same regressions with a
  \textbf{random intercept per frog}, so many cells from the same frog
  are not treated as independent repeats
\item
  \textbf{N/C vs cell size:} mixed model for log(N/C) vs log(cell area)
\end{itemize}
\end{keystats}

\section{Cell area distribution}\label{cell-area-distribution}

Distribution of cell areas (\(\mu\)m\(^2\)) for all good cells.

\begin{figure}[H]
\centering
\includegraphics[width=0.98\linewidth]{./figures/AreaDistribution.png}
\caption{Cell area distribution}
\end{figure}

\begin{keystats}
\begin{itemize}
\tightlist
\item
  \textbf{Median cell area: 290.69 \(\mu\)m\(^2\)} (IQR 257.70--335.48
  \(\mu\)m\(^2\)).
\item
  \textbf{Mean \(\pm\) SD: 303.10 \(\pm\) 67.07 \(\mu\)m\(^2\).}
\item
  \textbf{Circle equivalent: 19.24 \(\mu\)m} --- the diameter of a
  circle with the same area as the median cell (\(d = 2\sqrt{A/\pi}\)).
\end{itemize}
\end{keystats}

\section{Cell diameter distribution}\label{cell-diameter-distribution}

Long-axis and short-axis diameters (\(\mu\)m) for good cells. The two
distributions track each other closely.

\begin{figure}[H]
\centering
\includegraphics[width=0.98\linewidth]{./figures/DiameterDistribution.png}
\caption{Cell diameter distribution}
\end{figure}

\section{Nucleus size distribution}\label{nucleus-size-distribution}

\begin{figure}[H]
\centering
\includegraphics[width=0.98\linewidth]{./figures/NucleusDistribution.png}
\caption{Nucleus distribution}
\end{figure}

\section{Nucleus-to-cell area ratio
(N/C)}\label{nucleus-to-cell-area-ratio-nc}

Distribution of nucleus area divided by cell area.

\begin{figure}[H]
\centering
\includegraphics[width=0.98\linewidth]{./figures/NCRatioDistribution.png}
\caption{N/C ratio distribution}
\end{figure}

\begin{keystats}
\begin{itemize}
\tightlist
\item
  \textbf{Median N/C ratio: 0.13} (nucleus area \(\approx\) 0.13 of cell
  area --- about 13\% of the cell).
\item
  \textbf{Mean \(\pm\) SD: 0.14 \(\pm\) 0.04.}
\end{itemize}
\end{keystats}

\section{Per-frog variation in cell and nucleus
area}\label{per-frog-variation-in-cell-and-nucleus-area}

Boxplots for the top-40 frogs by good-cell count, sorted by per-frog
median cell area.

\begin{figure}[H]
\centering
\includegraphics[width=0.98\linewidth]{./figures/PerFrogBoxplot.png}
\caption{Per-frog cell area}
\end{figure}

\begin{figure}[H]
\centering
\includegraphics[width=0.98\linewidth]{./figures/PerFrogNucleusBoxplot.png}
\caption{Per-frog nucleus area}
\end{figure}

\begin{keystats}
Per-frog mean cell area spans \textbf{212.55--482.23 \(\mu\)m\(^2\)}.
The typical frog's average (median across frogs) was \textbf{290.34
\(\mu\)m\(^2\)}.
\end{keystats}

\section{Classification yield}\label{classification-yield}

Not every detected cell is measured in this report. After segmentation,
the quality classifier labels each cell \textbf{good}, \textbf{bad}, or
\textbf{rejected} (uncertain). \textbf{Yield} is the \textbf{good
fraction} --- the share of scored detections that were accepted as good
enough to measure. This section shows how that fraction varies across
frogs and images.

\textbf{What to expect in the figure below.} The \textbf{left panel} is
a stacked bar chart for the first 80 frogs: \textbf{green} = good cells,
\textbf{orange} = bad, \textbf{pink} = rejected. Taller green stacks
mean more cells passed the filter for that frog. The \textbf{right
panel} plots each frog's good fraction against its total good-cell count
--- it shows whether frogs with more accepted cells also tend to have
higher (or lower) yield. The summary numbers after the figure give
cohort-wide yield and the per-frog / per-image range.

\begin{figure}[H]
\centering
\includegraphics[width=0.98\linewidth]{./figures/YieldByFrog.png}
\caption{Classification yield}
\end{figure}

Of \textbf{383,915} scored detections, \textbf{57,520} were accepted as
good (\textbf{15.0\%}). Per-frog good fraction: median \textbf{17.7\%}
(range \textbf{2.7--62.1\%}). Per-image median good fraction:
\textbf{17.6\%}; \textbf{16} images had zero good cells.

\textbf{Important.} All cell-size results in this report use
\textbf{good cells only}. Cells labelled bad or rejected are
\textbf{not} included in area, diameter, or N/C summaries. When you read
those size numbers, also check \textbf{yield}: if a frog or image had a
\textbf{low good fraction}, fewer cells were measured there, so its size
average may be less reliable and the sample may not represent everything
that was seen in the microscope image.

\section{Nucleus--cell scaling
(regression)}\label{nucleuscell-scaling-regression}

This section asks: \textbf{when a cell gets larger, does the nucleus
grow in proportion?} We fit \textbf{log--log regressions} on good cells
with a matched nucleus (57,218 cells, 479 frogs). On a log--log plot, a
straight line means nucleus size scales as a \textbf{power} of cell
size. A slope of \textbf{1} is isometric scaling (nucleus and cell grow
at the same rate); \textbf{below 1} is negative allometry (the nucleus
grows more slowly than the cell).

\textbf{What to expect in the figures below.} There are \textbf{three
scatter plots}, one for each size pair:

\begin{enumerate}
\def\labelenumi{\arabic{enumi}.}
\tightlist
\item
  \textbf{Nucleus area vs cell area}
\item
  \textbf{Nucleus long axis vs cell long axis}
\item
  \textbf{Nucleus short axis vs cell short axis}
\end{enumerate}

Each plot uses \textbf{log--log axes} (both axes logarithmic). Every
point is one cell; the top frogs by cell count are shown in
\textbf{colour}, the rest in grey. Each plot overlays a \textbf{solid
black OLS line} (one best-fit line through all pooled cells) and a
\textbf{dashed orange Mixed line} (the within-frog trend with a random
intercept per frog).

\begin{figure}[H]
\centering
\includegraphics[width=0.98\linewidth]{./figures/NucleusVsCellRegression.png}
\caption{Nucleus area vs cell area}
\end{figure}

\begin{figure}[H]
\centering
\includegraphics[width=0.98\linewidth]{./figures/NucleusMajorVsCellMajor.png}
\caption{Nucleus major vs cell major axis}
\end{figure}

\begin{figure}[H]
\centering
\includegraphics[width=0.98\linewidth]{./figures/NucleusMinorVsCellMinor.png}
\caption{Nucleus minor vs cell minor axis}
\end{figure}

\textbf{What to expect in the table below.} One row per comparison
(area, long axis, short axis). The table gives the numeric slopes that
match the plots:

\begin{keystats}
\begin{itemize}
\tightlist
\item
  \textbf{OLS slope} --- ordinary least squares across \textbf{all
  cells} (each cell counted separately; mixes variation between and
  within frogs).
\item
  \textbf{Mixed slope} --- mixed-effects model with a \textbf{random
  intercept per frog}; estimates how nucleus size changes with cell size
  \textbf{among cells from the same frog}.
\item
  \textbf{OLS 95\% CI} --- uncertainty interval for the OLS slope.
\item
  \textbf{R\(^2\)} --- fraction of variance in log(nucleus) explained by
  log(cell) in the \textbf{OLS} model (does not apply to the mixed
  model).
\item
  \textbf{Slope \textless{} 1} in every row --- negative allometry:
  larger cells have proportionally smaller nuclei along that measure.
  Mixed slopes are \textbf{smaller} than OLS because much of the
  nucleus--cell correlation across the cohort is \textbf{between frogs},
  not within them.
\end{itemize}
\end{keystats}

\begin{longtable}[]{@{}
  >{\raggedright\arraybackslash}p{0.34\linewidth}
  >{\centering\arraybackslash}p{0.09\linewidth}
  >{\centering\arraybackslash}p{0.11\linewidth}
  >{\centering\arraybackslash}p{0.19\linewidth}
  >{\centering\arraybackslash}p{0.12\linewidth}
  >{\centering\arraybackslash}p{0.07\linewidth}@{}}
\caption{\textbf{Table --- Nucleus--cell scaling regressions} (good
cells with matched nucleus; 57,218 cells, 479 frogs). OLS = fit across
all cells; Mixed = random intercept per frog. R\(^2\) = fraction of
variance in log(nucleus) explained by log(cell) in the OLS
model.}\tabularnewline
\toprule\noalign{}
\begin{minipage}[b]{\linewidth}\raggedright
Comparison
\end{minipage} & \begin{minipage}[b]{\linewidth}\raggedright
n
\end{minipage} & \begin{minipage}[b]{\linewidth}\raggedright
OLS slope
\end{minipage} & \begin{minipage}[b]{\linewidth}\raggedright
OLS 95\% CI
\end{minipage} & \begin{minipage}[b]{\linewidth}\raggedright
Mixed slope
\end{minipage} & \begin{minipage}[b]{\linewidth}\raggedright
R\(^2\)
\end{minipage} \\
\midrule\noalign{}
\endfirsthead
\toprule\noalign{}
\begin{minipage}[b]{\linewidth}\raggedright
Comparison
\end{minipage} & \begin{minipage}[b]{\linewidth}\raggedright
n
\end{minipage} & \begin{minipage}[b]{\linewidth}\raggedright
OLS slope
\end{minipage} & \begin{minipage}[b]{\linewidth}\raggedright
OLS 95\% CI
\end{minipage} & \begin{minipage}[b]{\linewidth}\raggedright
Mixed slope
\end{minipage} & \begin{minipage}[b]{\linewidth}\raggedright
R\(^2\)
\end{minipage} \\
\midrule\noalign{}
\endhead
\bottomrule\noalign{}
\endlastfoot
Nucleus area vs cell area & 57,218 & 0.526 & {[}0.518, 0.533{]} & 0.149
& 0.255 \\
Nucleus long axis vs cell long axis & 57,218 & 0.616 & {[}0.609,
0.624{]} & 0.368 & 0.311 \\
Nucleus short axis vs cell short axis & 57,218 & 0.426 & {[}0.418,
0.434{]} & 0.182 & 0.148 \\
\end{longtable}

\section{N/C ratio vs cell size}\label{nc-ratio-vs-cell-size}

\textbf{N/C ratio vs cell size.} The plot below shows whether larger
cells have a proportionally smaller or larger nucleus (lower or higher
N/C). The table summarizes mixed-effects models on \textbf{log(N/C)
$\approx$ log(cell area)}.

\textbf{Random intercept per frog} --- each frog may have a different
baseline N/C; we test whether N/C still changes with cell area across
cells. \textbf{Random intercept + slope per frog} --- allows the
N/C--size trend to vary by frog as well. A \textbf{negative slope} means
larger cells tend to have \textbf{lower N/C} (a proportionally smaller
nucleus). Both models show a clear negative slope.

\begin{figure}[H]
\centering
\includegraphics[width=0.98\linewidth]{./figures/NCRatioVsArea.png}
\caption{N/C vs cell area}
\end{figure}

\begin{longtable}[]{@{}
  >{\raggedright\arraybackslash}p{(\linewidth - 8\tabcolsep) * \real{0.2000}}
  >{\raggedright\arraybackslash}p{(\linewidth - 8\tabcolsep) * \real{0.2000}}
  >{\raggedright\arraybackslash}p{(\linewidth - 8\tabcolsep) * \real{0.2000}}
  >{\raggedright\arraybackslash}p{(\linewidth - 8\tabcolsep) * \real{0.2000}}
  >{\raggedright\arraybackslash}p{(\linewidth - 8\tabcolsep) * \real{0.2000}}@{}}
\caption{\textbf{Table --- Mixed models for N/C vs cell area} (log--log;
good cells with matched nucleus). Slope = change in log(N/C) per unit
change in log(cell area).}\tabularnewline
\toprule\noalign{}
\begin{minipage}[b]{\linewidth}\raggedright
Model
\end{minipage} & \begin{minipage}[b]{\linewidth}\raggedright
Slope (log cell area)
\end{minipage} & \begin{minipage}[b]{\linewidth}\raggedright
95\% CI
\end{minipage} & \begin{minipage}[b]{\linewidth}\raggedright
n
\end{minipage} & \begin{minipage}[b]{\linewidth}\raggedright
Frogs
\end{minipage} \\
\midrule\noalign{}
\endfirsthead
\toprule\noalign{}
\begin{minipage}[b]{\linewidth}\raggedright
Model
\end{minipage} & \begin{minipage}[b]{\linewidth}\raggedright
Slope (log cell area)
\end{minipage} & \begin{minipage}[b]{\linewidth}\raggedright
95\% CI
\end{minipage} & \begin{minipage}[b]{\linewidth}\raggedright
n
\end{minipage} & \begin{minipage}[b]{\linewidth}\raggedright
Frogs
\end{minipage} \\
\midrule\noalign{}
\endhead
\bottomrule\noalign{}
\endlastfoot
Random intercept per frog & -0.8515 & {[}-0.8597, -0.8433{]} & 57,218 &
479 \\
Random intercept + slope per frog & -0.8134 & {[}-0.8342, -0.7926{]} &
57,218 & 479 \\
\end{longtable}

\section{Reference intervals}\label{reference-intervals}

Normative percentile ranges for good-cell morphology:

{\def\LTcaptype{none} % do not increment counter
\begin{scitable}
\begin{longtable}[]{@{}llll@{}}
\toprule\noalign{}
Metric & 2.5th \%ile & Median & 97.5th \%ile \\
\midrule\noalign{}
\endhead
\bottomrule\noalign{}
\endlastfoot
Cell area ($\mu$m$^2$) & 200.90 & 290.69 & 463.10 \\
Cell equivalent diameter ($\mu$m) & 15.99 & 19.24 & 24.28 \\
Cell long axis ($\mu$m) & 19.40 & 24.52 & 32.57 \\
Cell short axis ($\mu$m) & 12.15 & 15.23 & 19.36 \\
Nucleus area ($\mu$m$^2$) & 28.57 & 39.06 & 63.36 \\
N/C area ratio & 0.09 & 0.13 & 0.21 \\
Cell axis ratio & 1.25 & 1.62 & 2.08 \\
\end{longtable}
\end{scitable}
}

\section{Frog summary snapshot}\label{frog-summary-snapshot}

Ten frogs with the \textbf{smallest} mean cell area and ten with the
\textbf{largest}, plus a rank scatter of outlier frogs.

\begin{figure}[H]
\centering
\includegraphics[width=0.98\linewidth]{./figures/OutlierFrogs.png}
\caption{Outlier frogs}
\end{figure}

Ten frogs with the \textbf{smallest} mean cell area:

{\def\LTcaptype{none} % do not increment counter
\begin{scitable}
\begin{longtable}[]{@{}lllll@{}}
\toprule\noalign{}
frog\_id & n\_cells & area\_um2\_mean & nc\_ratio\_mean &
good\_fraction \\
\midrule\noalign{}
\endhead
\bottomrule\noalign{}
\endlastfoot
89.0 & 119.0 & 212.55 & 0.16 & 0.19 \\
41.0 & 187.0 & 216.02 & 0.19 & 0.18 \\
101.0 & 187.0 & 218.23 & 0.18 & 0.18 \\
35.0 & 138.0 & 219.15 & 0.14 & 0.21 \\
99.0 & 146.0 & 221.54 & 0.18 & 0.17 \\
73.0 & 114.0 & 221.88 & 0.17 & 0.12 \\
81.0 & 199.0 & 227.21 & 0.15 & 0.26 \\
30.0 & 127.0 & 227.88 & 0.16 & 0.24 \\
28.0 & 100.0 & 228.04 & 0.18 & 0.13 \\
250.0 & 168.0 & 229.04 & 0.15 & 0.21 \\
\end{longtable}
\end{scitable}
}

Ten frogs with the \textbf{largest} mean cell area:

{\def\LTcaptype{none} % do not increment counter
\begin{scitable}
\begin{longtable}[]{@{}lllll@{}}
\toprule\noalign{}
frog\_id & n\_cells & area\_um2\_mean & nc\_ratio\_mean &
good\_fraction \\
\midrule\noalign{}
\endhead
\bottomrule\noalign{}
\endlastfoot
221.0 & 89.0 & 449.54 & 0.11 & 0.09 \\
218.0 & 78.0 & 450.36 & 0.12 & 0.06 \\
210.0 & 116.0 & 454.62 & 0.12 & 0.23 \\
166.0 & 98.0 & 456.92 & 0.12 & 0.23 \\
431.0 & 99.0 & 457.59 & 0.16 & 0.13 \\
154.0 & 98.0 & 462.55 & 0.11 & 0.15 \\
455.0 & 83.0 & 465.68 & 0.11 & 0.48 \\
443.0 & 69.0 & 467.09 & 0.11 & 0.05 \\
119.0 & 90.0 & 469.63 & 0.11 & 0.08 \\
440.0 & 136.0 & 482.23 & 0.15 & 0.21 \\
\end{longtable}
\end{scitable}
}

\section{Discussion}\label{discussion}

This report describes the size and shape of \textbf{good-quality frog
blood cells} from a large automated microscopy cohort. Across
\textbf{57,520} accepted cells from \textbf{479} frogs, the typical cell
covered about \textbf{290.69 \(\mu\)m\(^2\)} (median area; middle 50\%
between \textbf{257.70} and \textbf{335.48 \(\mu\)m\(^2\)}), roughly the
area of a circle \textbf{19.24 \(\mu\)m} wide. That picture is useful as
a \textbf{reference description} of well-segmented, classifier-approved
cells in this dataset.

A striking feature is \textbf{how much average cell size differs between
frogs}. Per-frog mean areas ranged from \textbf{212.55} to
\textbf{482.23 \(\mu\)m\(^2\)}, while the typical frog's average (median
across frogs) was \textbf{290.34 \(\mu\)m\(^2\)}. Roughly \textbf{0.661}
of the total variation in cell area is associated with \textbf{which
frog} a cell came from; the rest reflects spread \textbf{within} each
frog (median within-frog CV about \textbf{11.7\%}). That pattern may
reflect real biological differences between frogs, differences in sample
preparation or imaging, or both; we cannot separate those sources here.
Frogs at the extremes of the size rank (see the frog summary snapshot
and outlier plot) are \textbf{candidates for manual review}: they may
represent genuine tail biology, technical artefacts, or low cell counts
that make the per-frog average less stable.

\textbf{Nucleus size relative to cell size} also carries a consistent
message. Nucleus--cell scaling is negatively allometric (nucleus grows
slower than cell area; larger cells are relatively less nucleus-dense).
In plain terms, when cells get larger, the nucleus tends to grow
\textbf{more slowly} than the cell body, so larger cells often show a
\textbf{lower N/C ratio} (a smaller nuclear share of cell area). That is
\textbf{consistent with} a model in which cytoplasm expands faster than
nuclear content during cell enlargement, though we have not tested
underlying mechanisms. The mixed-effects regressions suggest that part
of the nucleus--cell correlation across the cohort reflects
\textbf{between-frog} differences, not only within-frog
co-variation---so population-level trends should be read alongside
per-frog summaries.

\textbf{Classifier yield} shapes every morphometric statement in this
report. Only \textbf{15.0\%} of \textbf{383,915} scored detections were
labelled good (median per-frog good fraction \textbf{17.7\%}, range
\textbf{2.7--62.1\%}). Cells with uncertain quality were
\textbf{rejected} rather than forced into good or bad classes.
Morphometry therefore describes the \textbf{accepted subset}---cells the
model judged clear enough to measure---not the full pool of segmented
objects. Frogs or images with \textbf{low yield} contributed fewer
measured cells; \textbf{16} images had no good cells at all. When
interpreting a frog's mean area or N/C, it is worth checking its good
fraction in the exported summary tables: a sparse or selective accepted
set may not represent everything visible in the image.

\section{Appendix}\label{appendix}

Full per-frog metrics (all morphometry columns, yield, and flags) are
exported as \texttt{frog\_summary\_report.csv} and
\texttt{frog\_summary\_report.xlsx}.

\section{Glossary}\label{glossary}

{\def\LTcaptype{none} % do not increment counter
\begin{scitable}
\begin{longtable}[]{@{}
  >{\raggedright\arraybackslash}p{(\linewidth - 2\tabcolsep) * \real{0.4000}}
  >{\raggedright\arraybackslash}p{(\linewidth - 2\tabcolsep) * \real{0.6000}}@{}}
\toprule\noalign{}
\begin{minipage}[b]{\linewidth}\raggedright
Term
\end{minipage} & \begin{minipage}[b]{\linewidth}\raggedright
Meaning
\end{minipage} \\
\midrule\noalign{}
\endhead
\bottomrule\noalign{}
\endlastfoot
\textbf{Good cell} & A segmented cell accepted by the quality classifier
for morphometry. \\
\textbf{N/C ratio} & Nucleus area divided by cell area (unitless). \\
\textbf{Yield / good fraction} & Fraction of scored detections accepted
as good for a frog or image. \\
\textbf{ICC} & Share of cell-area variation explained by differences
between frogs (which frog a cell came from). \\
\textbf{Within-frog spread (CV)} & For each frog, (SD of cell areas
\(\div\) mean area) \(\times\) 100\%; report gives the median across
frogs. \\
\textbf{Typical frog's average} & Median of per-frog mean cell areas ---
each frog counts once, not weighted by cell count. \\
\textbf{Circle equivalent diameter} & Diameter of a circle with the same
area as the cell; a readable length summary, not the cell's true
width. \\
\textbf{Allometry} & How nucleus size scales with cell size; slope = 1
is isometric, \textless{} 1 is negative allometry. \\
\textbf{Reference interval} & Percentile range (e.g.~2.5th--97.5th)
describing typical good-cell morphology. \\
\end{longtable}
\end{scitable}
}

\end{document}
